% =========================================================
% Riemannian Geometry — Notes Index
% =========================================================

% ---------------------------------------------------------
% Breadcrumb
% ---------------------------------------------------------
\begin{tcolorbox}[
  colback=gray!6,
  colframe=gray!40,
  arc=2pt,
  left=8pt, right=8pt, top=6pt, bottom=6pt,
  title={\small\textbf{Where You Are in the Journey}},
  fonttitle=\small\bfseries
]
\begin{center}
\small
Differential Geometry
$\;\to\;$ \textbf{Riemannian Geometry}
$\;\to\;$ Stochastic Analysis on Manifolds
$\;\to\;$ Brownian Motion on a Torus
$\;\to\;$ $\cdots$
\end{center}

\medskip
\noindent\textbf{How we got here.}
Differential geometry showed how calculus extends to curves and
surfaces, and introduced the smooth manifold as the natural
domain.
Riemannian geometry equips each tangent space of a manifold
with an inner product that varies smoothly from point to point ---
the Riemannian metric.
This single addition restores all of the metric concepts
(distance, angle, volume, curvature) in the fully intrinsic setting.

\medskip
\noindent\textbf{What this chapter builds.}
Riemannian metrics and the metric tensor; the Levi-Civita
connection; geodesics as length-minimising curves and via the
geodesic equation; parallel transport; the Riemann curvature
tensor; sectional, Ricci, and scalar curvature; the
exponential map; and comparison geometry.

\medskip
\noindent\textbf{Where this leads.}
The Levi-Civita connection is the intrinsic derivative that
stochastic analysis on manifolds requires.
Hsu's \textit{Stochastic Analysis on Manifolds} builds
Brownian motion directly on Riemannian manifolds using the
connection, the exponential map, and the Laplace--Beltrami
operator --- all developed in this chapter.
\end{tcolorbox}
\vspace{1em}

% ---------------------------------------------------------
% Structural Role
% ---------------------------------------------------------
\begin{tcolorbox}[
  colback=gray!6,
  colframe=gray!40,
  arc=2pt,
  left=8pt, right=8pt, top=6pt, bottom=6pt,
  title={\small\textbf{Structural Role: Riemannian Geometry — Intrinsic Metric Structure}},
  fonttitle=\small\bfseries
]
Riemannian geometry is the culmination of the geometry sequence
and the direct prerequisite for the project's terminal goal:
stochastic analysis on manifolds.

The Riemannian metric is the manifold analogue of the inner product.
It makes precise what it means for two tangent vectors to be
perpendicular, for a curve to have a definite length, and for
a region to have a definite volume --- all without any reference
to an ambient space.

The Levi-Civita connection is the unique connection compatible
with the metric and torsion-free.
It defines parallel transport, covariant derivatives, and
the geodesic equation.
These are the tools that replace ordinary derivatives when
working on curved spaces, and they are exactly what is needed
to define stochastic differential equations on manifolds.

The Laplace--Beltrami operator, the Riemannian generalisation
of the ordinary Laplacian, is the infinitesimal generator of
Brownian motion on a Riemannian manifold.
\end{tcolorbox}
\vspace{1em}

% ---------------------------------------------------------
% Structural Roadmap
% ---------------------------------------------------------
\subsection*{Structural Roadmap}

\begin{center}
\textbf{Definitions $\longrightarrow$ Main Theorems
$\longrightarrow$ Consequences and Structural Insight}
\end{center}

The global progression is:

\begin{enumerate}
  \item Riemannian metrics: definition, examples, and the metric tensor
  \item Lengths of curves and the Riemannian distance function
  \item Connections and covariant derivatives
  \item The Levi-Civita connection: existence and uniqueness
  \item Parallel transport along curves
  \item Geodesics: the geodesic equation and the exponential map
  \item The Riemann curvature tensor: definition and symmetries
  \item Sectional curvature, Ricci curvature, and scalar curvature
  \item The Laplace--Beltrami operator
  \item Model spaces: spheres, hyperbolic space, flat tori
  \item Comparison geometry and the Hopf--Rinow theorem
\end{enumerate}

\begin{remark*}[Primary sources]
Lee, \textit{Riemannian Manifolds: An Introduction to Curvature};
do~Carmo, \textit{Riemannian Geometry};
Hsu, \textit{Stochastic Analysis on Manifolds}, Chapter~1
(for the stochastic-analysis perspective on the Riemannian toolkit).
\end{remark*}

% ---------------------------------------------------------
% Status
% ---------------------------------------------------------
\vspace{1em}
\begin{tcolorbox}[
  colback=gray!6, colframe=gray!40, arc=2pt,
  left=8pt, right=8pt, top=6pt, bottom=6pt,
  title={\small\textbf{Status: Planned}},
  fonttitle=\small\bfseries
]
\begin{center}\Large\bfseries Coming Soon\end{center}
\vspace{6pt}
\noindent Notes, proofs, and exercises will appear here in a future revision.
\end{tcolorbox}

% Future sections (uncomment as content is written):
% \input{volume-v/geometry/riemannian-geometry/notes/notes-metric-tensor}
% \input{volume-v/geometry/riemannian-geometry/notes/notes-levi-civita}
% \input{volume-v/geometry/riemannian-geometry/notes/notes-geodesics}
% \input{volume-v/geometry/riemannian-geometry/notes/notes-curvature-tensor}
% \input{volume-v/geometry/riemannian-geometry/notes/notes-laplace-beltrami}
% \input{volume-v/geometry/riemannian-geometry/notes/notes-model-spaces}
